%%%%%%%%%%  scan idx 456  |  folio 445  %%%%%%%%%%

\medskip
\noindent
10.5. \underline{Reducción de 10.4 a un caso particular}

\smallskip
\noindent
10.5.1. La unicidad de un $u$ que satisfaga la condición enunciada
(incluso para un solo $\ell$) es clara. Por otra parte, la existencia
significa simplemente que cada $u_\ell$ aplica de hecho
$\underline M\otimes\underline M'$ en $\mathbb Z_S$, y que el
homomorfismo así obtenido
\begin{equation*}
\underline M\otimes\underline M'\longrightarrow\mathbb Z_S
\end{equation*}
es independiente de $\ell$. De hecho, basta incluso probar la
afirmación análoga para los
$u_\ell\otimes_{\mathbb Z_\ell}\mathbb Q_\ell$ y para un
homomorfismo
\begin{equation*}
u_{\mathbb Q}:\underline M_{\mathbb Q}\otimes
 \underline M'_{\mathbb Q}\longrightarrow\mathbb Q_S.
\end{equation*}

Supongamos entonces establecida la existencia después de un cambio de
base $S'\longrightarrow S$ como en 9.3. Se deduce inmediatamente de
(10.3.5) que existe un homomorfismo
\begin{equation*}
u:\underline M\otimes\underline M'
 \longrightarrow\frac{1}{e(S'/S)}\mathbb Z_S\subset\sisi{Q_S}{\mathbb Q_S}
\end{equation*}
compatible con los
$u_\ell\otimes_{\mathbb Z_\ell}\mathbb Q_\ell$; esto prueba a). En
cuanto a b), es entonces claro que, si vale para $A_{K'}$, también vale
para $A_K$.

Por consiguiente, al probar 10.4 podemos efectuar cualquier cambio
dominante de rasgos $S'\longrightarrow S$. Esto permite, en particular,
suponer que $S$ es completo y tiene cuerpo residual algebraicamente
cerrado. En este caso, $\underline M$ y $\underline M'$ son esquemas en
grupos constantes, de valores $M$ y $M'$, módulos libres de tipo finito
sobre $\mathbb Z$, y a) y b) se reformulan ventajosamente en términos
de estos últimos.

\smallskip
\noindent
10.5.2. Supongamos ahora que $A_K$ es isomorfo a un factor directo de
un esquema abeliano $\overline A_K$ sobre $K$ con reducción
semiestable, y que 10.4 ha sido probado para este último. Se deduce
enseguida que 10.4 vale también para $A_K$, habida cuenta del evidente
comportamiento funtorial de los retículos $\underline M,\underline M'$
y de los acoplamientos de monodromía respecto de las sumas

%%%%%%%%%%  scan idx 457  |  folio 446  %%%%%%%%%%

directas de esquemas abelianos con reducción semiestable.

Es igualmente inmediato, gracias a 10.5.1, que la validez de 10.4 solo
depende de la clase de isogenia de $A_K$. Junto con la observación
precedente y el ``teorema de reducibilidad de Poincaré'' [15], esto
muestra que, si $A_K$ es cociente de un esquema abeliano
$\overline A_K$ para el cual vale 10.4, entonces 10.4 también vale para
$A_K$.

\smallskip
\noindent
10.5.3. Utilicemos ahora el hecho bien conocido [31, teorema 11
(p.~20)] de que, tras efectuar si es necesario una extensión finita de
$K$ (lo cual está permitido por 10.5.1), todo esquema abeliano $A_K$ es
cociente de la jacobiana
\begin{equation*}
\underline{\operatorname{Pic}}^0_{X_K/K}
\end{equation*}
de una curva propia, lisa y geométricamente conexa $X_K$ sobre
\sisi{$X$}{$K$}. En virtud de 10.5.2, quedamos reducidos a probar 10.4
para una jacobiana de este tipo.

\smallskip
\noindent
10.5.4. Según un teorema de Artin--Mumford [3] (ya citado y utilizado
en el Exposé V), tras efectuar si es necesario una nueva extensión
finita $K'$ de $K$, podemos suponer que $X_K$ se prolonga a un esquema
proyectivo y plano sobre $S$, regular y cuya fibra especial geométrica
$X_{\bar s}$ no tiene más singularidades que singularidades cuadráticas
ordinarias (es decir, en este caso, cruces normales de dos ramas). En
este caso probaremos 10.4 a) en 12.5, describiendo por completo el
acoplamiento de monodromía (10.4.1) en términos de los datos geométricos
de la situación; al mismo tiempo verificaremos 10.4 b) para la
polarización canónica de la jacobiana.

%%%%%%%%%%  scan idx 458  |  folio 447  %%%%%%%%%%

\smallskip
\noindent
10.5.5. Queda por ver que, si para un $A_K$ dado con reducción
semiestable vale la afirmación 10.4 a), y la afirmación 10.4 b) vale
para una polarización dada $\xi_o$ de $A_K$, entonces la misma
conclusión vale para toda polarización $\xi$ de $A_K$. En efecto, sean
\begin{equation*}
\widetilde\xi_o,\widetilde\xi:A_K\longrightarrow A'_K
\end{equation*}
los homomorfismos definidos por $\xi_o$ y $\xi$; estos son isogenias.
Entonces
\begin{equation*}
\widetilde\xi=\widetilde\xi_o\alpha
 ={}^t(\alpha^*)\widetilde\xi_o,
 \qquad
\alpha\in\operatorname{End}(A_K)\otimes_{\mathbb Z}\mathbb Q,
\end{equation*}
donde $\alpha\longmapsto\alpha^*$ denota la involución canónica del
álgebra $\operatorname{End}(A_K)\otimes_{\mathbb Z}\mathbb Q$
definida por la polarización $\xi_o$, y el exponente $t$ en
${}^t(\alpha^*)$ denota la traspuesta (que asocia a un endomorfismo de
$A_K$ un endomorfismo de $A'_K$). Por otra parte, se sabe [15,
cap.~V, teorema 3] que
\begin{equation*}
\alpha=\sum_i\beta_i^*\beta_i,
 \qquad
\beta_i\in\operatorname{End}(A_K)\otimes_{\mathbb Z}\mathbb Q,
\end{equation*}
y, por tanto,
\begin{equation*}
\widetilde\xi=\sum_i{}^t\beta_i\,\widetilde\xi_o\,\beta_i.
\end{equation*}
Equivalentemente, en términos de las biextensiones (es decir, de las
correspondencias divisoriales) asociadas a las polarizaciones
consideradas (VIII 4.14),
\begin{equation}\tag{10.5.5.1}
\delta(\xi)=\sum_i\delta(\xi_o)\circ(\beta_i,\beta_i).
\end{equation}
Según lo dicho después de (10.2.7), esta fórmula implica (con la
notación de (10.4.2))
\begin{equation}\tag{10.5.5.2}
\varphi^t_\xi
 =\sum_i\varphi^t_{\xi_o}(\beta_i,\beta_i)
 \in\operatorname{Hom}
 \bigl(M_{\mathbb Q}\otimes_{\mathbb Q}M_{\mathbb Q},\mathbb Q\bigr).
\end{equation}

%%%%%%%%%%  scan idx 459  |  folio 448  %%%%%%%%%%

Se deduce, en efecto, que, si $\varphi^t_{\xi_o}$ es simétrica y
positiva, también lo es $\varphi^t_\xi$.

Finalmente, la no degeneración de $\varphi^t_{\xi_o}$ equivale
evidentemente a la afirmación de que el acoplamiento (10.4.1) es
separante (pues
$\xi_{ot}:M\otimes_{\mathbb Z}\mathbb Q
\longrightarrow M'\otimes_{\mathbb Z}\mathbb Q$ es un isomorfismo, ya
que $\widetilde\xi_o$ es una isogenia), y por tanto es asimismo
independiente de la polarización elegida. También cabe observar que,
para verificar que (10.4.1) es separante, basta verificar la misma
afirmación para el acoplamiento de monodromía $u_\ell$ de (9.1.2) para
un número primo $\ell$. Tomando $\ell\ne p$, se puede utilizar la
descripción de $u_\ell$ dada en 9.2; el hecho de que $u_\ell$ sea
separante por la izquierda (y, por tanto, por ambos lados, por razones
de rango) es entonces inmediato.
